\documentclass[a4paper,12pt]{ctexart}
\usepackage{xcolor,graphicx,amsmath}
\usepackage[top=1.8cm, bottom=1.8cm, left=1.8cm, right=1.8cm]{geometry}
\usepackage{array,hyperref}
\hypersetup{colorlinks=true,linkcolor=blue!70!black}\linespread{1.35}
\definecolor{defblue}{RGB}{0,80,160}\definecolor{thinkorange}{RGB}{230,120,20}
\definecolor{funboxbg}{RGB}{245,248,252}\definecolor{practicegreen}{RGB}{40,130,70}
\newenvironment{thinkbox}{\vspace{0.3cm}\begin{center}\begin{minipage}{0.88\textwidth}\colorbox{funboxbg}{\begin{minipage}{0.94\textwidth}\vspace{0.15cm}\textbf{\textcolor{thinkorange}{【 想一想 】}}}{\vspace{0.15cm}\end{minipage}}\end{minipage}\end{center}\vspace{0.3cm}}
\newenvironment{definition}[1]{\vspace{0.2cm}\begin{center}\begin{minipage}{0.88\textwidth}\fbox{\begin{minipage}{0.92\textwidth}\vspace{0.1cm}\textbf{\textcolor{defblue}{【 #1 】}}}{\vspace{0.1cm}\end{minipage}}\end{minipage}\end{center}\vspace{0.2cm}}
\newenvironment{practice}{\vspace{0.3cm}\begin{center}\begin{minipage}{0.88\textwidth}\colorbox{funboxbg}{\begin{minipage}{0.94\textwidth}\vspace{0.15cm}\textbf{\textcolor{practicegreen}{【 练一练 】}}}{\vspace{0.15cm}\end{minipage}}\end{minipage}\end{center}\vspace{0.2cm}}
\newenvironment{figplaceholder}[1]{\vspace{0.2cm}\begin{center}\fbox{\begin{minipage}{0.82\textwidth}\centering\textcolor{blue!70!black}{\textbf{【插图位置】}}\\[0.2cm]#1\end{minipage}}\end{center}\vspace{0.2cm}}
\title{\textcolor{blue!70!black}{\Huge\textbf{地理2 · 地貌、灾害与资源环境}}}
\author{}\date{}
\begin{document}\maketitle\thispagestyle{empty}
\section*{\textcolor{blue!70!black}{致同学}}
地理1讲了地球的"基本设定"——自转公转、气候、水循环。这一本讲地球表面的\textbf{造型师}——是什么力量塑造了高山和峡谷？自然灾害是怎么回事？我们的淡水和能源从哪来？以及——地理学家用什么工具研究这些。

\newpage
\section{\textcolor{orange!80!black}{第一课\quad 海陆分布与板块构造}}

\subsection*{七大洲四大洋}
亚洲（最大）、非洲、北美洲、南美洲、南极洲、欧洲、大洋洲。\\
太平洋（最大）、大西洋、印度洋、北冰洋（最小）。

\subsection*{板块构造学说（复习与深化）}
地球岩石圈分为六大板块：亚欧板块、太平洋板块、印度板块、非洲板块、美洲板块、南极洲板块。

\textbf{板块边界的地质作用：}
\begin{itemize}
  \item 大洋中脊（板块张裂）：岩浆上涌→新洋壳形成（如大西洋中脊）。
  \item 俯冲带（板块碰撞）：大洋板块俯冲入大陆板块之下→海沟+岛弧/火山链（如环太平洋火山地震带——日本、印尼）。
  \item 大陆碰撞带：两个大陆板块相撞→山脉隆起（印度板块撞上亚欧板块→喜马拉雅山）。
\end{itemize}

\begin{thinkbox}
珠穆朗玛峰每年"长高"约1厘米——不是它自己在长，而是印度板块还在继续挤压亚欧板块。北京在地质上属于华北板块（亚欧板块的一部分）——在漫长的地质历史上，华北还曾是一片海洋。
\end{thinkbox}

\begin{practice}
1. 地球岩石圈分为\_\_\_\_\_大板块。\\
2. 喜马拉雅山的形成是由于\_\_\_\_\_板块和\_\_\_\_\_板块碰撞。\\
3. 环太平洋火山地震带位于板块的\_\_\_\_\_（边界/内部）。\\
4. 四大洋中最大的是\_\_\_\_\_；最小的是\_\_\_\_\_。\\
5. （判断）板块内部地壳比较稳定。（\quad）\\
6. 环太平洋地震带被称为"\_\_\_\_\_"。\\
7. 大西洋中脊位于板块的\_\_\_\_\_（张裂/碰撞）边界。\\
8. 华北板块属于\_\_\_\_\_板块的一部分。
\end{practice}

\newpage
\section{\textcolor{orange!80!black}{第二课\quad 地形地貌——内外力作用}}

\subsection*{内力作用——地球内部的"动力"}

\textbf{褶皱}：岩层在挤压下弯曲——形成山脉。向斜（凹槽→储水）和背斜（拱起→储油/气）。\\
\textbf{断层}：岩层断裂错动——形成地垒（如华山/庐山）和地堑（如渭河平原/汾河谷地）。\\
\textbf{火山}：岩浆沿地壳薄弱处喷出。

\subsection*{外力作用——太阳和水在"削平"地球}
\textbf{风化}：温度变化、水的冻融、生物作用——把岩石"弄碎"。\\
\textbf{侵蚀（主要类型）：}
\begin{itemize}
  \item 河流侵蚀：V形谷（上游下切）、U形谷（下游侧蚀）。桂林山水——流水的化学溶蚀（喀斯特地貌）。
  \item 风力侵蚀：风蚀蘑菇、雅丹地貌。中国的魔鬼城就是典型的风蚀地貌。
  \item 冰川侵蚀：U形谷、角峰、刃脊。西藏很多湖泊和山谷都是冰川"挖"出来的。
  \item 海浪侵蚀：海蚀崖、海蚀柱。
\end{itemize}
\textbf{搬运：}风、水、冰川把碎屑物带走。\\
\textbf{沉积：}碎屑物堆积——形成三角洲（河口）、冲积扇（山前）、沙丘（沙漠）。

\begin{thinkbox}
内力和外力是一对"对头"——内力把地球表面弄得"坑坑洼洼"（造山），外力"努力把它磨平"（风化和侵蚀）。两者较量了几十亿年，形成了我们今天看到的千姿百态的地表。
\end{thinkbox}

\begin{practice}
1. 背斜是岩层向\_\_\_\_\_弯曲，向斜是岩层向\_\_\_\_\_弯曲。\\
2. 桂林山水（喀斯特地貌）是由\_\_\_\_\_水的\textbf{\_\_\_\_\_}（溶蚀/沉积）形成的。\\
3. （判断）内力作用只会使地表变得平坦。（\quad）\\
4. 黄河三角洲是\_\_\_\_\_作用形成的。\\
5. 褶皱中向斜是良好的\_\_\_\_\_（储水/储油）构造。\\
6. 风蚀蘑菇是\_\_\_\_\_（风力侵蚀/流水侵蚀）形成的。\\
7. （判断）华山属于断层中的地垒构造。（\quad）\\
8. 河流上游以下切侵蚀为主，常形成\_\_\_\_\_形谷。
\end{practice}

\newpage
\section{\textcolor{orange!80!black}{第三课\quad 土壤与植被}}

\subsection*{土壤——地球的"皮肤"}
土壤由矿物质（来自岩石风化）、有机物（来自动植物残体分解）、空气和水分组成。土壤剖面从上到下：表土层（腐殖质含量高、暗色）→ 心土层 → 底土层（母质层）。

\textbf{北京的主要土壤类型：}褐土和潮土。褐土在山区和丘陵（偏碱性），潮土在平原区的河流两岸。

\subsection*{植被与气候的对应}
\textbf{森林}：降水充足（年降水量$>400\ \text{mm}$）。北京周边的百花山就属于温带落叶阔叶林。\\
\textbf{草原}：降水中等（$200-400\ \text{mm}$）。内蒙古高原。\\
\textbf{荒漠}：降水极少（$<200\ \text{mm}$）。新疆的塔克拉玛干沙漠。

\begin{practice}
1. 土壤由\_\_\_\_\_、\_\_\_\_\_、\_\_\_\_\_和\_\_\_\_\_四部分组成。\\
2. 年降水量200-400mm的地区，天然植被多为\_\_\_\_\_。\\
3. （判断）森林植被只分布在热带地区。（\quad）\\
4. 北京平原地区的土壤类型主要是\_\_\_\_\_。\\
5. 土壤剖面中腐殖质含量最高的是\_\_\_\_\_层。\\
6. 年降水量大于\_\_\_\_\_mm的地区一般发育森林植被。\\
7. （判断）荒漠植被区年降水量小于200mm。（\quad）\\
8. 土壤中\_\_\_\_\_来自岩石风化，\_\_\_\_\_来自动植物残体分解。
\end{practice}

\newpage
\section{\textcolor{orange!80!black}{第四课\quad 自然灾害}}

\subsection*{地质灾害}
\textbf{地震}（在自然故事中已详述——此处复习要点）：板块运动→岩石断裂→地震波。震级越大能量越大（每增加1级——能量约增32倍）。烈度表示地面受破坏程度——同一地震烈度各地不同（离震中越远烈度越小）。

\textbf{滑坡与泥石流}：在山区——暴雨或地震后——大量土石沿山坡滑下或与水混合形成泥石流。防护：建排水沟、加固坡面、植树造林。北京西部的房山/门头沟山区是泥石流易发区。

\textbf{火山}：岩浆喷出地表。活火山（正在或可能喷发）、休眠火山（历史上有但长时间未喷发）、死火山（不可能再喷发）。长白山天池是一座休眠火山。

\subsection*{气象灾害}
\textbf{台风}（热带气旋）：在西太平洋生成——登陆中国东南沿海。带来强风+暴雨+风暴潮。北京不受台风直接影响（离海较远），但台风外围常常给华北带来大量水汽——引发暴雨。

\textbf{洪涝}：降雨过多或排水不畅。北京"七下八上"（7月下旬到8月上旬）是主汛期——历史上多次发生城市内涝。

\textbf{旱灾}：长时间降水偏少——农业减产、人畜饮水困难。华北地区经常春夏连旱——南水北调的重要目的之一就是缓解华北旱情。

\begin{figplaceholder}
此处插入中国主要自然灾害分布图（见 figure/requirements/geo2-ch1.txt 插图1）
\end{figplaceholder}

\begin{practice}
1. 地震震级每增加1级——能量约增\_\_\_\_\_倍。\\
2. 台风带来\_\_\_\_\_、\_\_\_\_\_和\_\_\_\_\_。\\
3. 北京的主汛期是\_\_\_\_\_。\\
4. （判断）泥石流多发生在坡度较缓的河谷地区。（\quad）\\
5. 火山按活动状态分为\_\_\_\_\_、\_\_\_\_\_和\_\_\_\_\_。\\
6. 长白山天池目前属于\_\_\_\_\_火山。\\
7. （判断）地震震级越大，烈度一定越大。（\quad）\\
8. 华北地区经常发生春旱，\_\_\_\_\_是缓解旱情的重要工程。
\end{practice}

\newpage
\section{\textcolor{orange!80!black}{第五课\quad 自然资源}}

\subsection*{水资源}
地球淡水仅占总水量的约2.5\%——且绝大部分冻在南北极冰川中。中国是缺水国家——人均水资源仅为世界平均水平的1/4。

\textbf{南水北调：}把长江流域的水调往华北（北京的饮用水很大一部分来自丹江口水库）。这是人类史上最大规模的跨流域调水工程。

\subsection*{土地资源}
中国"人多地少"——用全球7\%的耕地养活了全球20\%的人口。耕地主要集中在东部季风区（华北平原、长江中下游平原、东北平原）。\textbf{18亿亩耕地红线}——确保粮食安全的底线。

\subsection*{矿产资源与能源}
\textbf{煤炭}：中国最重要的化石能源（山西、陕西、内蒙古三省份占全国产量约70\%）。但燃煤排放CO$_2$和污染物——需要大力发展清洁能源。\\
\textbf{石油和天然气}：中国进口量大。\\
\textbf{可再生能源}：太阳能（西北——甘肃省酒泉的光伏基地）、风能（内蒙古/沿海——北京延庆冬季风大）、水能（长江三峡）。

\subsection*{中国的能源转型}
"双碳"目标：2030年前碳达峰、2060年前碳中和。减少化石燃料——增加非化石能源（核能、水电、风电、光伏）的比例。

\begin{practice}
1. 地球上淡水占总水量的比例约\_\_\_\_\%。\\
2. 中国的耕地红线是\_\_\_\_\_亿亩。\\
3. 北京饮用水的重要来源之一是\_\_\_\_\_。\\
4. （判断）中国是石油出口国。（\quad）\\
5. 中国能源消费中占比最大的是\_\_\_\_\_。\\
6. 中国太阳能资源最丰富的地区是\_\_\_\_\_。\\
7. （判断）我国耕地主要集中在西部干旱区。（\quad）\\
8. 碳中和是指\_\_\_\_\_的排放量和\_\_\_\_\_量相抵消。
\end{practice}

\newpage
\section{\textcolor{orange!80!black}{第六课\quad 地理技术与北京在地探索}}

\subsection*{遥感（RS）——从太空看地球}
卫星上的传感器接收地表反射的电磁波——获得影像。利用遥感可以监测：北京的城市扩张（从1970年代到今天的卫星对比）、森林火灾监测、沙尘暴路径追踪。

\subsection*{GIS——地理信息系统}
把地理信息"数字化"——叠加不同图层（地图、人口、交通、土地利用）——进行分析和决策。疫情追踪地图、北京公交路线规划——背后都是GIS。

\subsection*{GPS/北斗——定位与导航}
北斗卫星导航系统——中国自主建设的全球卫星定位系统。精度可达米级。出租车调度、手机地图导航——都离不开定位系统。

\subsection*{北京在地项目探索}（与拓展3呼应）
自然地理视角的一些"在家门口"可以做的事情：
\begin{itemize}
  \item \textbf{西山山地考察}：观察褶皱和断层、识别岩石（北京西山的石灰岩、花岗岩等）。
  \item \textbf{永定河生态修复}：之前的沙石采挖导致河道干涸——近年通过补水修复。
  \item \textbf{城市热岛效应}：北京城区和郊区的温差——自己用温度计记录对比。
  \item \textbf{密云水库水源地调查}：了解北京的水源来自哪里、周边生态保护措施。
\end{itemize}

\begin{thinkbox}
地理不只是"背地名"——它是理解地球这个复杂系统运行方式的一门学科。你学会了自然地理（气候怎么来的、山怎么形成的）——又结合了生物（生态系统）、物理（热力学/天体运动）、化学（岩石风化/水污染）——你已经拥有了理解整个**地球系统**的知识框架。
\end{thinkbox}

\begin{practice}
1. RS的中文名是\_\_\_\_\_；GIS的中文名是\_\_\_\_\_。\\
2. 中国的自主卫星导航系统叫\_\_\_\_\_。\\
3. （判断）城市热岛指的是城市气温高于郊区的现象。（\quad）\\
4. 永定河是北京的\_\_\_\_\_（重要河流/不重要）。\\
5. 卫星导航系统可以实现\_\_\_\_\_级定位精度。\\
6. （判断）GIS可以叠加不同图层进行分析。（\quad）\\
7. 遥感卫星可以监测北京的城市\_\_\_\_\_。\\
8. 北京西山的地质考察可以观察到\_\_\_\_\_和\_\_\_\_\_等地质构造。
\end{practice}

\vspace{0.5cm}
\begin{center}\rule{0.7\textwidth}{1pt}\vspace{0.4cm}
{\large\textcolor{blue!70!black}{\textbf{—— 地理系列（2本）完 ——}}}
\end{center}

\newpage
\section*{\textcolor{defblue}{参考答案}}

\subsection*{第一课}
1. 六\\
2. 印度、亚欧\\
3. 边界\\
4. 太平洋、北冰洋\\
5. （√）\\
6. 环太平洋火山地震带\\
7. 张裂\\
8. 亚欧

\subsection*{第二课}
1. 上（拱起）、下（凹陷）\\
2. 流、溶蚀\\
3. （×）内力作用使地表变得高低不平\\
4. 沉积\\
5. 储水\\
6. 风力侵蚀\\
7. （√）\\
8. V

\subsection*{第三课}
1. 矿物质、有机物、空气、水分\\
2. 草原\\
3. （×）温带、寒带都有森林\\
4. 潮土\\
5. 表土（腐殖质层）\\
6. 400\\
7. （√）\\
8. 矿物质、有机物

\subsection*{第四课}
1. 约32\\
2. 强风、暴雨、风暴潮\\
3. "七下八上"（7月下旬\~{}8月上旬）\\
4. （×）多发生在坡度较陡的山区\\
5. 活火山、休眠火山、死火山\\
6. 休眠\\
7. （×）烈度还与距离震中远近等有关\\
8. 南水北调

\subsection*{第五课}
1. 2.5\%\\
2. 18\\
3. 丹江口水库（南水北调）\\
4. （×）中国是石油进口国\\
5. 煤炭\\
6. 西北（如甘肃/新疆/青海）\\
7. （×）集中在东部季风区\\
8. CO$_2$、吸收（或固定）

\subsection*{第六课}
1. 遥感、地理信息系统\\
2. 北斗\\
3. （√）\\
4. 重要河流\\
5. 米\\
6. （√）\\
7. 扩张（变化）\\
8. 褶皱、断层

\end{document}
